% Part: first-order-logic
% Chapter: completeness
% Section: introduction

\documentclass[../../../include/open-logic-section]{subfiles}

\begin{document}

\iftag{FOL}
      {\olfileid{fol}{com}{int}}
      {\olfileid{pl}{com}{int}}

\olsection{تعارف}

تمامیت دا قضیہ منطق دے سب توں بنیادی نتائج وچوں اک اے۔ ایس دے
دو بیان نیں، جنہاں دی ہم ارزیت اسیں ثابت کراں گے۔ پہلے بیان
مطابق ایہ معنوی نتیجے تے ساڈے !!{derivation} نظام دے تعلق بارے
بنیادی گل کہندا اے: جے کوئی !!a{sentence}~$!A$ کجھ
!!{sentence}s~$\Gamma$ دا نتیجہ ہووے، تاں اک
!!a{derivation} وی موجود اے جیہڑا $\Gamma
\Proves !A$ قائم کردا اے۔ ایس طرح !!{derivation} نظام
اُنا طاقت ور اے جِنا اوہ انہاں گلاں نوں ثابت کیتے بنا ہو سکدا اے
جیہڑیاں حقیقت وچ نتیجہ نہیں بندیاں۔

دوجے بیان مطابق ایہ نمونے دے وجود دا نتیجہ اے: !!{sentence}s
دا ہر بے تضاد سیٹ پورا ہون جوگ اے۔ بے تضادی اک ثبوتی-نظریاتی
تصور اے: ایہ کہندا اے کہ ساڈا !!{derivation} نظام کسے خاص
طرح دیاں !!{derivation}s نہیں بنا سکدا۔ پر صرف ایس گل توں کہ
~$\Gamma$ توں ایس طرح دیاں !!{derivation}s نہیں نکلدیاں، ایہ
کیویں پکا ہویا کہ
% PNB-COM-012: FOL شاخ وچ وی دتے مفروضیاں نوں پورا کرن والی ساخت دی شرط لازمی اے۔
\iftag{FOL}{!!a{structure}~$\Struct{M}$ موجود ہووے تے دتے سیٹ نوں پورا کرے}{!!{valuation}{~$\pAssign{v}$}
موجود ہووے جیہدی لئی $\iftag{FOL}{\Sat{M}}{\pSat{v}}{\Gamma}$}؟
تمامیت دا قضیہ پہلی واری ثابت ہون توں پہلاں---بلکہ ساڈے ہُن
والے !!{derivation} نظاماں توں وی پہلاں---جرمنی دے عظیم ریاضی
دان ڈیوڈ ہلبرٹ دی رائے سی کہ ریاضیاتی نظرییاں دی بے تضادی
انہاں دے موضوع بنن والیاں چیزاں دے وجود دی ضمانت دیندی اے۔
انہاں نے گوٹلوب فریگے نوں اک خط وچ ایویں لکھیا:
\begin{quote}
  جے من مانی طرح دتے مسلّمے، اپنے سارے نتیجیاں سمیت، اک دوجے
  نال تضاد نہ کردے ہون، تاں اوہ سچے نیں تے مسلّماں نال تعریف
  کیتیاں چیزاں موجود نیں۔ میرے لئی سچائی تے وجود دا معیار ایہی اے۔
\end{quote}
فریگے نے ایس نال سخت اختلاف کیتا۔ تمامیت دے قضیے دا دوجا
بیان دکھاندا اے کہ ہلبرٹ گھٹ توں گھٹ ایس معنی وچ درست سن کہ
جے مسلّمے بے تضاد ہون، تاں \emph{کوئی نہ کوئی}
\iftag{FOL}{!!{structure}}{!!{valuation}} موجود اے جیہڑی
انہاں ساریاں نوں سچا بناندی اے۔

ایہو وجوہاں تمامیت دے قضیے---یا خاص طور تے اوس دے ثبوت---دی
اہمیت دیاں اکلوتیاں وجوہاں نہیں۔ ایس دے کئی اہم نتائج نیں،
جنہاں وچوں کجھ اُتے اسیں الگ گل کراں گے۔ مثال لئی، کوئی وی
!!{derivation} جیہڑا $\Gamma \Proves !A$ وکھاندا اے محدود
ہوندا اے، ایس لئی~$\Gamma$ دے صرف محدود گنتی دے
!!{sentence}s ورت سکدا اے۔ اس توں تمامیت دے قضیے نال ایہ
نکلدا اے کہ جے $!A$،~$\Gamma$ دا نتیجہ اے، تاں اوہ پہلے ای
~$\Gamma$ دے کسے محدود ذیلی سیٹ دا نتیجہ اے۔ ایس نوں
\emph{محدود گواہی (compactness)} آکھدے نیں۔ برابر دا بیان
ایہ اے کہ جے $\Gamma$ دا ہر محدود ذیلی سیٹ بے تضاد اے، تاں
$\Gamma$ وی بے تضاد اے۔

بھاویں محدود گواہی دا قضیہ، تمامیت دے قضیے توں
!!{derivation}s راہیں گھوم کے نکلدا اے، پر تمامیت دے قضیے دے
\emph{ثبوت} نوں ورت کے ایس نوں سِدھا ثابت کرنا وی ممکن اے۔
ثبوت اک ایسی خاصیت---بے تضادی---والا !!{sentence}s دا سیٹ
لیندا اے تے ایس سیٹ توں !!a{structure} بناندا اے جیہدیاں
کجھ خاصیتاں ہوندیاں نیں (ایتھے ایہ کہ اوہ سیٹ نوں پورا کردی اے)۔
تقریباً ایہو تعمیر، بے تضاد سیٹاں دی تھاں ”محدود طور تے پورے
ہون جوگ“ !!{sentence}s دے سیٹاں توں شروع کر کے، محدود گواہی
سِدھی قائم کر سکدی اے۔
\iftag{FOL}{ایس تعمیر دے ہور نتیجے وی نیں؛ مثال لئی،
  !!{sentence}s دے ہر پورے ہون جوگ سیٹ دا اک محدود یا
  !!{denumerable}s نمونہ موجود اے۔ (ایس نتیجے نوں
  L\"owenheim--Skolem دا قضیہ آکھدے نیں۔) عام طور تے
  !!{sentence}s دے سیٹاں توں !!{structure}s بنان دی
  تعمیر منطق وچ اکثر ورتی جاندی اے، تے کدی فلسفے وچ وی۔}{}

\end{document}
